\documentclass[10pt,a4paper]{amsart}
\usepackage{amsmath,amssymb,amsthm,graphicx}
\usepackage{verbatim}
\newenvironment{reference}{\comment}{\endcomment}
%\newenvironment{reference}{}{}
\newenvironment{slogan}{\comment}{\endcomment}
\newenvironment{history}{\comment}{\endcomment}
% For commutative diagrams we use Xy-pic
\usepackage[all]{xy}

% We use 2cell for 2-commutative diagrams.
\xyoption{2cell}
% Theorem environments.
%
\theoremstyle{plain}
\newtheorem{theorem}[subsection]{Theorem}
\newtheorem{proposition}[subsection]{Proposition}
\newtheorem{lemma}[subsection]{Lemma}

\theoremstyle{definition}
\newtheorem{definition}[subsection]{Definition}
\newtheorem{example}[subsection]{Example}
\newtheorem{exercise}[subsection]{Exercise}
\newtheorem{situation}[subsection]{Situation}

\theoremstyle{remark}
\newtheorem{remark}[subsection]{Remark}
\newtheorem{remarks}[subsection]{Remarks}

\numberwithin{equation}{subsection}

% Macros
%
\def\lim{\mathop{\mathrm{lim}}\nolimits}
\def\colim{\mathop{\mathrm{colim}}\nolimits}
\def\Spec{\mathop{\mathrm{Spec}}}
\def\Hom{\mathop{\mathrm{Hom}}\nolimits}
\def\Ext{\mathop{\mathrm{Ext}}\nolimits}
\def\SheafHom{\mathop{\mathcal{H}\!\mathit{om}}\nolimits}
\def\SheafExt{\mathop{\mathcal{E}\!\mathit{xt}}\nolimits}
\def\Sch{\mathit{Sch}}
\def\Mor{\mathop{\mathrm{Mor}}\nolimits}
\def\Ob{\mathop{\mathrm{Ob}}\nolimits}
\def\Sh{\mathop{\mathit{Sh}}\nolimits}
\def\NL{\mathop{N\!L}\nolimits}
\def\CH{\mathop{\mathrm{CH}}\nolimits}
\def\proetale{{pro\text{-}\acute{e}tale}}
\def\etale{{\acute{e}tale}}
\def\QCoh{\mathit{QCoh}}
\def\Ker{\mathop{\mathrm{Ker}}}
\def\Im{\mathop{\mathrm{Im}}}
\def\Coker{\mathop{\mathrm{Coker}}}
\def\Coim{\mathop{\mathrm{Coim}}}

% Boxtimes
%
\DeclareMathSymbol{\boxtimes}{\mathbin}{AMSa}{"02}

%
% Macros for moduli stacks/spaces
%
\def\QCohstack{\mathcal{QC}\!\mathit{oh}}
\def\Cohstack{\mathcal{C}\!\mathit{oh}}
\def\Spacesstack{\mathcal{S}\!\mathit{paces}}
\def\Quotfunctor{\mathrm{Quot}}
\def\Hilbfunctor{\mathrm{Hilb}}
\def\Curvesstack{\mathcal{C}\!\mathit{urves}}
\def\Polarizedstack{\mathcal{P}\!\mathit{olarized}}
\def\Complexesstack{\mathcal{C}\!\mathit{omplexes}}
% \Pic is the operator that assigns to X its picard group, usage \Pic(X)
% \Picardstack_{X/B} denotes the Picard stack of X over B
% \Picardfunctor_{X/B} denotes the Picard functor of X over B
\def\Pic{\mathop{\mathrm{Pic}}\nolimits}
\def\Picardstack{\mathcal{P}\!\mathit{ic}}
\def\Picardfunctor{\mathrm{Pic}}
\def\Deformationcategory{\mathcal{D}\!\mathit{ef}}
\usepackage{cleveref}
\crefname{originalreference}{original reference}{original references}
\Crefname{originalreference}{Original reference}{Original references}
\makeatletter
\newlabel{topology-definition-generic-point}{{topology-definition-generic-point}{source}{}{sourcepointer-1}{}}
\newlabel{topology-definition-generic-point@cref}{{[originalreference][0][]topology-definition-generic-point}{[1][0][]source}{}{}{}}
\makeatother
\crefname{originalreference}{original reference}{original references}
\Crefname{originalreference}{Original reference}{Original references}
\makeatletter
\newlabel{topology-lemma-Noetherian}{{topology-lemma-Noetherian}{source}{}{sourcepointer-2}{}}
\newlabel{topology-lemma-Noetherian@cref}{{[originalreference][0][]topology-lemma-Noetherian}{[1][0][]source}{}{}{}}
\makeatother
\title{The space of a locally Noetherian algebraic stack is quasi-sober}
\author{The Stacks Project Authors}
\date{stacks-project:stacks-morphisms.tex:local-snapshot-20261012; GNU FDL 1.2 or later}
\begin{document}
\maketitle
\noindent\textit{Editorial note.} All delivered mathematical text is the author's own, copied from the pinned Stacks Project chapter source sources/stacks-algebraic-stacks-morphisms/pinned-original-source.tex; no statement or proof step was rewritten, reordered or removed.\par\medskip
\noindent\textit{Editorial note.} Preamble provenance: the Stacks Project keeps one single shared preamble for every chapter; this card takes it directly from the staged sources/stacks-algebraic-stacks-morphisms/preamble.tex as three verbatim spans (lines 8-12, 14-18 and 154-217), so no derived file was written and nothing was added or redefined. The wrapper class and the base packages are supplied by the assembler.\par\medskip
\noindent\textit{Editorial note.} Bibliography provenance: neither the counted statement nor its proof contains a citation (this lane rescans the whole span before writing the card), so the card carries no bibliography and no entry was reconstructed.\par\medskip
\noindent\textit{Editorial note.} Selection provenance: this card was opened under the file-based control contract (CONTROL.md), with the mechanical pre-card gate run on its statement and proof interval before the card was written.\par\medskip
\noindent\textit{Allowed standard prerequisites (editorial).}\begin{itemize}
\item The notion of a quasi-sober topological space: every irreducible closed subset has a generic point. Invoked in the statement (source lines 2282-2283) as a pointer to the chapter on topology. Fix what has to be produced for an irreducible closed subset.
\item The finiteness of the number of irreducible components of a closed subset of a Noetherian space, and the openness of the map on topological spaces induced by a smooth morphism. Both are used inside the counted proof (source lines 2290-2296) as cross-chapter pointers and from the theory of presentations. Decomposes the preimage of the closed subset into finitely many components and shows that the image of one of them is dense.
\end{itemize}\medskip
\ifdefined\hypertarget\hypertarget{sourcepointer-1}{}\fi
\noindent\textit{Source reference.} Original reference topology-definition-generic-point: sources/stacks-algebraic-stacks-morphisms/pinned-original-source.tex, source lines 2283--2283. Navigation pointer only: the counted text refers to the Stacks item topology-definition-generic-point, which belongs to another chapter of the project and is not part of this card. The pointer carries the target's own label as its number because this chapter's staged snapshot does not contain the other chapter's tag table; its content is not reproduced here.\par\medskip
\ifdefined\hypertarget\hypertarget{sourcepointer-2}{}\fi
\noindent\textit{Source reference.} Original reference topology-lemma-Noetherian: sources/stacks-algebraic-stacks-morphisms/pinned-original-source.tex, source lines 2292--2292. Navigation pointer only: the counted text refers to the Stacks item topology-lemma-Noetherian, which belongs to another chapter of the project and is not part of this card. The pointer carries the target's own label as its number because this chapter's staged snapshot does not contain the other chapter's tag table; its content is not reproduced here.\par\medskip
\begin{lemma}
\label{lemma-locally-Noetherian-quasi-sober}
Let $\mathcal{X}$ be a locally Noetherian algebraic stack.
Then $|\mathcal{X}|$ is quasi-sober (Topology, Definition
\ref{topology-definition-generic-point}).
\end{lemma}

\begin{proof}
We have to prove that every
irreducible closed subset $T \subset |\mathcal{X}|$ has a generic point.
Choose an affine scheme $U$ and a smooth morphism $f : U \to \mathcal{X}$
such that $f^{-1}(T) \subset |U|$ is nonempty. Since $U$ is Noetherian,
the closed subset $f^{-1}(T)$ has finitely many irreducible components
(Topology, Lemma \ref{topology-lemma-Noetherian}).
Say $f^{-1}(T) = Z_1 \cup \ldots \cup Z_n$ is the decomposition into
irreducible components. As $f$ is open, the image of
$f|_{f^{-1}(T)} : f^{-1}(T) \to T$ contains a nonempty open subset of $T$.
Since $T$ is irreducible, this means that $f(f^{-1}(T))$ is dense.
Since $T$ is irreducible, it follows that $f(Z_i)$ is dense for some $i$.
Then if $\xi_i \in Z_i$ is the generic point we see that
$f(\xi_i)$ is a generic point of $T$.
\end{proof}

\end{document}
